\section{Study Design}
\label{sec:method}

\paragraph{Tasks and generator.}
\begin{itemize}
\item Prior-study tasks, games, functional tests; target ApoMario.
\item \emph{Highscore}: persist score, player name, survival time;
sorted menu display; 16 functional tests.
\item \emph{Achievements}: gameplay milestones; 20,000 points per run,
3 enemies killed, 60 seconds survived; persistence and menu display;
28 functional tests; prior study's hardest task.
\item \emph{Generation}: ApoMario only. \emph{Reuse}: also ApoIcarus,
with existing feature implementation.
\item Model: OpenAI \texttt{gpt-5.6-luna}, reasoning high, Codex adapter;
same model for acquisition, generation, judging.
\item Cell: task variant + base context; structural S, functional F,
behavioural B, alone/combined/none. Strategy comparison: Generation S
and Reuse S+B; 5 artifacts/cell, $N=10$ per arm.
\end{itemize}

\paragraph{Delivery.}
\begin{itemize}
\item \emph{One response}: no feedback; strict delivery---reject
whole-file replacement of existing files and edits with missing anchors.
Baseline also under lenient delivery: whole-file replacement accepted
(Table~\ref{tab:matrix}), matching prior pipeline.
\item \emph{Agentic}: repository search/read tools; compiler and test
feedback; at most 5 submissions within 24 tool turns.
\item \emph{Judge}: separate conversation; pending submission + context;
cancellation with quoted submitted line violating a cited requirement/control;
advice returned; at most 3 cancellations/trajectory.
\end{itemize}

\paragraph{Security context.}
\begin{itemize}
\item Acquisition: one-line task + vocabulary/schema; read-only repository
where applicable; shell inspection under command budget.
\item \emph{S1}: repository-derived high-level guidance, no code citations.
\item \emph{S2}: full document, statements anchored to code; citations
resolved against parsed source; statements without resolvable evidence dropped.
\item \emph{S3}: generic document, no repository access.
\item Document appended once after task; acquisition instructions v13
throughout strategy tables; explicit numeric bounds and validation domains;
earlier revisions/effects in study repository.
\end{itemize}

\begin{table}[t]
\centering
\footnotesize
\setlength{\tabcolsep}{4pt}
\caption{Ten security checks: categories, weakness types, properties, requirement sources. CWE~\citep{cwe}; feature contract; OWASP ASVS~\citep{asvs5}; CERT Oracle Secure Coding Standard for Java~\citep{long2011}.}
\label{tab:mapping}
\resizebox{\textwidth}{!}{\begin{tabular}{@{}llll@{}}
\toprule
Category & Weakness type & Property & Requirement source \\
\midrule
Input policy & CWE-20, CWE-1284 & Integrity & Contract; ASVS validation; CERT IDS \\
Retention policy & CWE-770 & Availability & Contract \\
Parser robustness & CWE-755, CWE-20 & Availability & CERT ERR \\
Resource stress & CWE-400, CWE-770, CWE-789 & Availability & CERT MSC, FIO \\
Deserialization dispatch & CWE-502 & Integrity & CERT SER; ASVS \\
\bottomrule
\end{tabular}}
\end{table}

\paragraph{Evaluation.}
\begin{itemize}
\item Isolated JVM; functional tests + 10 scored security checks/task
(Table~\ref{tab:mapping}); unrunnable check = not passed.
\item Functional tests: Highscore, 7 unit + 4 integration + 5 live-game;
Achievements, 16 + 7 + 5.
\item Highscore input policy (5): reject negative score/time,
missing/blank/over-long name. Retention (1): at most 100 records.
\item Achievements input policy (5): reject negative points/time;
prevent point-total overflow; reject unknown/over-long stored IDs.
Retention (1): deduplicate stored IDs.
\item Both tasks (4): corrupted-file recovery; survive 64~MiB line
and million-record/line store; no Java object deserialization.
\item ``All pass'': all functional tests + 10 scored checks + 11th
persistence round-trip check; valid record/unlock survives restart.
\item Denominators: Compiled/All pass, all $N$ artifacts;
Tests/Checks, mean pass share among compiled artifacts.
\item Cost: generator + judge input tokens and model-call minutes/artifact.
\end{itemize}
